% !TeX root = ../regression_logistique.tex
\chapter{Construction du coût par vraisemblance}
\label{ann:cout}

\orient{%
écrire la vraisemblance d'un jeu d'étiquettes sous le modèle ; la transformer en
somme par le logarithme ; en déduire la perte employée par le fil principal ;
nommer les hypothèses qui interviennent}{%
logarithme, exposants (\annexeref{ann:prerequis}) ; chapitre~\ref{ch:erreur}}{%
approfondissement ; le fil principal n'emploie que la formule encadrée}

\section{Probabilité d'une réponse observée}

Notons $y_i$ l'étiquette : $1$ si l'observation $i$ appartient à \Gryf, $0$ sinon, et
$p_i=\sigma(z_i)$ la probabilité produite par le modèle. La probabilité que celui-ci
attribue à l'étiquette effectivement observée vaut $p_i$ si $y_i=1$, et $1-p_i$
si $y_i=0$. Ces deux cas se rassemblent :
\[
  \Pp(y_i)=p_i^{\,y_i}\,(1-p_i)^{\,1-y_i}.
\]

Vérification des deux cas, en se souvenant que $a^{0}=1$ et $a^{1}=a$ :
\[
  y_i=1:\quad p_i^{1}(1-p_i)^{0}=p_i\cdot1=p_i,
  \qquad
  y_i=0:\quad p_i^{0}(1-p_i)^{1}=1\cdot(1-p_i)=1-p_i .
\]
L'exposant nul ramène à $1$ le facteur qui ne correspond pas au cas observé.

\section{Vraisemblance du fichier}

Une seconde hypothèse intervient ici, distincte de celle du lien logit : les
étiquettes sont supposées conditionnellement indépendantes, une fois les notes
connues. C'est elle qui autorise à écrire la probabilité d'observer l'ensemble
des étiquettes du fichier comme le produit des probabilités individuelles :
\[
  L(\w)=\prod_{i=1}^{m}p_i^{\,y_i}(1-p_i)^{\,1-y_i}.
\]
$L$ s'appelle la \emph{vraisemblance}. Elle dépend des coefficients $\w$, car ce
sont eux qui produisent les $p_i$. Chercher les meilleurs coefficients, c'est
chercher ceux qui rendent $L$ maximale.

\begin{vigilance}[Sous-dépassement du produit]
Le produit compte $1\,600$ facteurs, tous strictement inférieurs à $1$, et
décroît donc très vite. Le plus petit double normal vaut environ
$2{,}2\times10^{-308}$ ; en dessous, la précision se dégrade puis le nombre
s'arrondit à zéro~\cite{goldberg1991}.

Le seuil se franchit facilement. Si chaque facteur valait $0{,}5$ --- un modèle
sans opinion --- le produit vaudrait $0{,}5^{1600}\approx10^{-482}$, qui
s'arrondit à zéro. Si chacun valait $0{,}9$, le produit vaudrait
$6{,}1\times10^{-74}$ : représentable, mais déjà privé de toute marge, et un
fichier dix fois plus long le ferait sous-dépasser. Une fois deux vraisemblances
arrondies à zéro, aucune comparaison n'est plus possible.

Le logarithme supprime le problème en transformant le produit en somme, dont
l'ordre de grandeur reste celui du nombre d'observations.
\end{vigilance}

\begin{figure}[htbp]
\centering
\begin{graphique}{Chaîne de trois expressions reliées par des flèches : le produit des probabilités individuelles, puis sa transformation en somme par le logarithme, puis le passage à une moyenne changée de signe, quantité que le programme minimise.}
\begin{tikzpicture}[
  bx/.style={draw=ink!30,rounded corners=2.5pt,align=center,font=\footnotesize,
    inner sep=5pt,text width=26mm,minimum height=13mm},
  op/.style={font=\scriptsize,text=accent,align=center},
  arr/.style={-{Latex[length=2.2mm]},thick,ink!60}]
\node[bx,fill=accentlight] (p) {$m$ probabilités\\[1pt] \footnotesize
  $\Pp(y_i)\in\,]0,1[$};
\node[bx,fill=warninglight,right=12mm of p] (l) {$L(\w)$\\[1pt] \footnotesize
  $10^{-482}$ à $\Mzero$};
\node[bx,fill=rulelight,right=12mm of l] (ll) {$\log L(\w)$\\[1pt] \footnotesize
  somme de $m$ termes};
\node[bx,fill=accentlight,right=12mm of ll] (j) {$J(\w)$\\[1pt] \footnotesize
  $0{,}693$ au départ};
\draw[arr] (p) -- node[op,above=1mm]{produit} (l);
\draw[arr] (l) -- node[op,above=1mm]{$\log$} (ll);
\draw[arr] (ll) -- node[op,above=1mm]{$\times\,(-1/m)$} (j);
\node[font=\scriptsize,text=ink!75,below=1.5mm of l,text width=34mm,
  align=center] {inutilisable en machine};
\node[font=\scriptsize,text=ink!75,below=1.5mm of j,text width=34mm,
  align=center] {à minimiser};
\end{tikzpicture}
\end{graphique}
\caption{Le produit des probabilités individuelles mesure la plausibilité du jeu de
coefficients. Le logarithme le convertit en somme à optimum inchangé ; le
facteur $-1/m$ change la maximisation en minimisation et rend la valeur
indépendante de l'effectif.}
\label{fig:vraisemblance}
\end{figure}
\FloatBarrier

\section{Passage au logarithme}

La vraisemblance est un produit de mille six cents facteurs, forme peu commode à
dériver et impossible à évaluer en machine. Le logarithme la transforme en somme
sans déplacer le point où elle est maximale, puisqu'il est strictement croissant.

\begin{align}
  \log L(\w)
  &= \log\prod_{i=1}^{m}p_i^{\,y_i}(1-p_i)^{\,1-y_i} \notag\\[2pt]
  &= \sum_{i=1}^{m}\log\Bigl[p_i^{\,y_i}(1-p_i)^{\,1-y_i}\Bigr]
    &&\text{le log d'un produit, sur les } m \text{ facteurs} \notag\\[2pt]
  &= \sum_{i=1}^{m}\Bigl[\log p_i^{\,y_i}+\log(1-p_i)^{\,1-y_i}\Bigr]
    &&\text{le log d'un produit, sur deux termes} \notag\\[2pt]
  &= \sum_{i=1}^{m}\Bigl[y_i\log p_i+(1-y_i)\log(1-p_i)\Bigr]
    &&\log a^{b}=b\log a \label{eq:loglik}
\end{align}

Le passage au logarithme est légitime pour chercher un maximum : $\log$ est
strictement croissante, donc elle conserve l'ordre. Le $\w$ qui maximise $L$ est
exactement celui qui maximise $\log L$.

\section{Changement de signe et moyenne}

Les algorithmes d'optimisation minimisent par convention. On change donc le
signe, et on divise par $m$ pour obtenir une valeur moyenne, indépendante de
l'effectif :
\begin{equation}
  \boxed{\;J(\w)=-\frac{1}{m}\sum_{i=1}^{m}
  \Bigl[y_i\log p_i+(1-y_i)\log(1-p_i)\Bigr]\;}
  \tag{$\dagger$}\label{eq:cout}
\end{equation}
Maximiser $\log L$ et minimiser $J$ sont la même opération : multiplier par
$-1/m$ retourne la fonction et la met à l'échelle, sans déplacer l'argument où
l'optimum est atteint. C'est la formule que donne le sujet du projet.

\section{Forme du terme individuel}

Le terme d'une observation vaut $-\log p_i$ si son étiquette est \Gryf, et
$-\log(1-p_i)$ sinon : dans les deux cas, l'opposé du logarithme de la
probabilité accordée à l'étiquette observée.

\begin{table}[htbp]
\centering
\begin{tabular}{ccrr}
\toprule
$y_i$ & $p_i$ & probabilité accordée à l'étiquette & perte $-\log(\cdot)$ \\
\midrule
$1$ & $0{,}99$ & $0{,}99$ & $0{,}010$ \\
$1$ & $0{,}75$ & $0{,}75$ & $0{,}288$ \\
$1$ & $0{,}50$ & $0{,}50$ & $0{,}693$ \\
$1$ & $0{,}10$ & $0{,}10$ & $2{,}303$ \\
$1$ & $0{,}01$ & $0{,}01$ & $4{,}605$ \\
\midrule
$0$ & $0{,}10$ & $0{,}90$ & $0{,}105$ \\
$0$ & $0{,}90$ & $0{,}10$ & $2{,}303$ \\
\bottomrule
\end{tabular}
\caption{Perte selon la probabilité accordée à l'étiquette observée : presque nulle pour
une prédiction correcte et assurée, divergente quand cette probabilité tend vers
zéro.}
\label{tab:cout-individuel}
\end{table}

\begin{figure}[htbp]
\centering
\begin{graphique}{Courbe de la perte en fonction de la probabilité accordée à l'étiquette observée : proche de zéro lorsque cette probabilité est proche de 1, divergente lorsqu'elle tend vers zéro.}
\begin{tikzpicture}
\begin{axis}[courseplot,height=6cm,axis lines=left,
  xlabel={probabilité accordée à la classe observée},ylabel={coût},
  xmin=0,xmax=1,ymin=0,ymax=5,samples=250,domain=.008:1]
  \addplot[accent,very thick]{-ln(x)};
  \addplot[only marks,mark=*,warning,mark size=2.4pt]
    coordinates {(0.99,0.010) (0.5,0.693) (0.1,2.303)};
  \node[warning,anchor=south west,font=\small] at (axis cs:.5,.75){$\log 2=0{,}693$};
\end{axis}
\end{tikzpicture}
\end{graphique}
\caption{La perte d'une prédiction en fonction de la probabilité que le modèle
attribuait à la classe observée. Le point $0{,}693$ correspond à un modèle qui
produit $0{,}5$.}
\label{fig:cout-individuel}
\end{figure}
\FloatBarrier

La formule rend compte du contrôle employé depuis le chapitre~\ref{ch:erreur} :
des coefficients nuls donnent $p_i=\sigma(0)=0{,}5$ pour chaque observation, donc
une perte moyenne de $-\log(0{,}5)=\log 2$.
